\section{Intrinsic dimension, lattice and growth}

A body of zero volume in $\mathbb R^n$ can still detect growth of lower order. The group generated by the valuation semigroup determines a degree period and a lattice in which the body is full-dimensional, and counting in this lattice gives the growth exponent and the leading coefficient. Ratios of sections identify the exponent with a transcendence degree, and since the coefficient is a limit of graded dimensions, it does not depend on the flag.

Throughout this section, some positive-degree piece of $W_\bullet$ is nonzero. Let $G\subset\mathbb Z^{n+1}$ be the subgroup generated by $S$, and let $\operatorname{pr}_{n+1}$ be the projection to the last coordinate. Define the \emph{degree period} $e\geq1$ by
\[
 \operatorname{pr}_{n+1}(G)=e\mathbb Z.
\]
Thus $e$ is the greatest common divisor of the degrees $m$ with $W_m\neq0$. Set
\[
 \Lambda=G\cap(\mathbb R^n\times\{0\}),
 \qquad d=\operatorname{rank}\Lambda.
\]
We regard $\Lambda$ as a lattice in its real span, a $d$-dimensional subspace of $\mathbb R^n\times\{0\}$. Its lattice measure assigns volume one to a fundamental parallelepiped. On each translate of this span we use the translated measure, denoted by $\operatorname{vol}_\Lambda$; when $d=0$, a point has volume one.

\begin{theorem}\label{thm:intrinsic-flag-count}\index{Convex body!intrinsic lattice volume}\index{Degree period}
The affine dimension of $\Delta(W_\bullet)$ is $d$. Its lattice volume is finite and strictly positive, and
\[
 \lim_{q\to\infty}\frac{\dim W_{eq}}{(eq)^d}
 =\operatorname{vol}_\Lambda\Delta(W_\bullet).
\]
If $e$ does not divide $m$, then $W_m=0$. In particular, $W_m\neq0$ for all sufficiently large multiples $m$ of $e$.
\end{theorem}
\begin{proof}
Choose $g_0=(b,e)\in G$ and an integral basis $\lambda_1,\ldots,\lambda_d$ of $\Lambda$. Every element of $G$ has a unique expression
\[
 \sum_{i=1}^d a_i\lambda_i+qg_0,
 \qquad (a,q)\in\mathbb Z^d\times\mathbb Z.
\]
These coordinates identify $G$ with $\mathbb Z^{d+1}$. Let $S'$ be the image of $S$ under this identification. It generates $\mathbb Z^{d+1}$, and the height of the image of $(x,m)\in S$ is $q=m/e\geq0$.

The points $(m^{-1}x,1)$, for $(x,m)\in S$ with $m>0$, span the $(d+1)$-dimensional real vector space $\operatorname{span}_{\mathbb R}G$. Since they all have height one, their affine span has dimension $d$. Taking the convex hull and the closure does not change this affine span, which proves the dimension assertion.

Intersect the cone containing $S$ with $\operatorname{span}_{\mathbb R}G$. The linear change of coordinates carries this intersection to a closed cone with compact height-one slice whose only point at height zero is the origin. Indeed, a vector $(x,m)$ with $|x|\leq Cm$ has horizontal coordinates bounded by a constant times $m$, hence by a constant times $q$. The height-one body of $S'$ is
\[
 \Delta'=e\bigl(\Delta(W_\bullet)-b/e\bigr),
\]
where the right-hand side is expressed in the basis $\lambda_1,\ldots,\lambda_d$. It is full-dimensional in $\mathbb R^d$, and therefore
\[
 \operatorname{vol}_{\mathbb R^d}\Delta'
 =e^d\operatorname{vol}_\Lambda\Delta(W_\bullet)>0.
\]
Apply the semigroup counting theorem of \cref{sec:graded-flag-body} to $S'$, which generates its full lattice. The points of $S'$ at height $q$ correspond bijectively to the valuations of the nonzero elements of $W_{eq}$, and by \cref{lem:flag-one-dimensional-leaves} there are $\dim W_{eq}$ of them. Hence
\[
 \lim_{q\to\infty}\frac{\dim W_{eq}}{q^d}
 =\operatorname{vol}_{\mathbb R^d}\Delta'.
\]
Dividing by $e^d$ gives the asserted limit, and its positivity gives eventual nonvanishing. For $d=0$ the body is a point of volume one, so the integer dimensions converge to one and are eventually equal to one. Finally, every degree with a nonzero section lies in the projection $e\mathbb Z$ of $G$, which proves the vanishing outside these multiples.
\end{proof}

\Needspace{6\baselineskip}
The growth dimension also has a description independent of the flag. Write $\mathcal M(X)$ for the field of meromorphic functions on $X$ and define the \emph{section-ratio field}
\[
 K(W_\bullet)=\{s/t:s,t\in W_m\text{ for some }m\geq0,
                                  \ t\neq0\}\subset\mathcal M(X).
\]
It is a field: sums and products of two such ratios, with denominators $t_1\in W_{m_1}$ and $t_2\in W_{m_2}$, can be written with the common denominator $t_1t_2$ in degree $m_1+m_2$, and inversion exchanges numerator and denominator. It contains $\mathbb C$ because $W_0=\mathbb C$.

\begin{proposition}\label{prop:flag-dimension-section-field}\label{prop:full-flag-body-meromorphic}\index{Section-ratio field}
The transcendence degree of $K(W_\bullet)$ over $\mathbb C$ is exactly $d$. If $\Delta(W_\bullet)$ has nonempty interior in $\mathbb R^n$, then $X$ has $n$ algebraically independent meromorphic functions. Moreover, $d$, the period $e$ and the intrinsic volume in \cref{thm:intrinsic-flag-count} are the same for any two complete flags on $X$.
\end{proposition}
\begin{proof}
Suppose first that $d>0$. Choose sections $s_0,\ldots,s_d$ of positive degrees $m_0,\ldots,m_d$ whose normalized valuations $\nu(s_i)/m_i$ are affinely independent. Let $M$ be a common multiple of their degrees and put $t_i=s_i^{M/m_i}\in W_M$. The ratios $f_i=t_i/t_0$ belong to $K(W_\bullet)$. The flag valuation extends to quotients by subtraction, giving
\[
 \nu(f_i)=M\left(\frac{\nu(s_i)}{m_i}
                  -\frac{\nu(s_0)}{m_0}\right),
 \qquad 1\leq i\leq d.
\]
These vectors are linearly independent. Distinct monomials in the $f_i$ therefore have distinct valuations, so the term of smallest valuation in a nonzero polynomial expression cannot cancel, and the ratios are algebraically independent. Thus the transcendence degree is at least $d$; this is trivial when $d=0$.

Conversely, let $f_1,\ldots,f_r$, with $r\geq1$, be algebraically independent elements of the field. Using the product of their denominators, write $f_i=A_i/B$ with $A_i,B\in W_M$ and $B\neq0$ for a common positive degree $M$. For every integer $q\geq1$ and every multi-index $a=(a_1,\ldots,a_r)\in\mathbb N^r$ with $|a|=a_1+\cdots+a_r\leq q$, the sections
\[
 B^{q-|a|}A_1^{a_1}\cdots A_r^{a_r}\in W_{Mq}
\]
are linearly independent, because division by $B^q$ turns a linear relation among them into a polynomial relation among the $f_i$. There are $\binom{q+r}{r}$ such sections. Since $W_M\neq0$, the degree $M$ is a multiple of $e$, and \cref{thm:intrinsic-flag-count} gives $\dim W_{Mq}=O(q^d)$. Comparing with the binomial lower bound forces $r\leq d$, which proves equality of the transcendence degree with $d$. A body with nonempty interior has $d=n$, and this gives the stated consequence for $X$.

The field and the set of degrees with nonzero sections do not depend on the flag, so neither $d$ nor $e$ does. By \cref{thm:intrinsic-flag-count}, the lattice volume is, for either flag, the limit of the same dimensions with the same normalization. It is therefore independent of the flag as well, although the bodies and their lattices may differ.
\end{proof}
